% ************************** Thesis Abstract *****************************
% Use `abstract' as an option in the document class to print only the titlepage and the abstract.
\begin{abstract}
Patch-based diffusion models promise to enable the generation of strong image priors in situations with limited training data and memory. This work reproduces and extends the Patch Diffusion Inverse Solver (PaDIS) prior and sampling methods presente by Hu et al. for computed tomography (CT), using LIDC-IDRI data and a fully specified, noise-free fan-beam geometry implemented in LION. To separate ambiguities in the original PaDIS paper from adaptations required by the new geometry, we develop three implementations: Paper-Faithful, matching the original PaDIS paper depiction; Repo-Faithful, matching the original PaDIS code; and LION-Physical, using power iteration to determine Lipschitz constants for data-consistency normalisation. These are evaluated on CT reconstruction with 8-, 20-, and 60-views, limited-angles and both $256\times 256$ and $512\times512$ resolutions, along with alternative samplers, different patch construction methods, and training ablations.

PaDIS produced high-quality reconstructions, including $40.08\,\mathrm{dB}$ PSNR and $0.960$ SSIM at $512\times512$. With the LION-Physical Lipschitz scaling, its separately tuned VE-DDNM configuration achieved the best PSNR at 20 and 60 views, reaching $36.49$ and $41.55\,\mathrm{dB}$. However, we did not reproduce the original PaDIS claim that a PaDIS prior combined with VE-DPS consistently outperforms whole-image diffusion. We found that a whole-image model performed better at both 20 and 8 views when each prior used separately tuned hyperparameters. We also found that the largest tested patches gave the strongest patch-size result under shared inference setting, positional encoding had little effect on conditioned reconstruction, and training on all available CT slices did not improve PaDIS under a fixed training-time budget. These results show that patch priors are viable and memory-scalable, but their advantage is conditional on the sampler, data regime, and treatment of the forward operator. Reproduction was further limited by incomplete geometry, training, tuning, and implementation details in the original PaDIS paper and by material differences between its algorithms and released code. The complete project codebase is available \href{https://github.com/THartigan/MPhil-Project}{here}, and the training runs and saved models are available \href{https://wandb.ai/tjh200-university-of-cambridge/PaDIS-Reproduction?nw=nwusertjh200}{here}.
\end{abstract}
